\exercisesection{Duality of modules of finite length}

\begin{enumerate}
\def\labelenumi{\arabic{enumi})}
\tightlist
\item
  Let A be a Noetherian local ring, and M an A-module. Show that for M to be Artinian, it is necessary and sufficient that its socle S be of finite dimension over \(\kappa_{\Lambda}\) and that M be an essential extension of S (that is, by (A, II, p.~185, exerc. 15), every nonzero submodule of M contains a nonzero element of S).
\end{enumerate}

¶ 2) Let A be a complete local Noetherian ring, and let M be a Λ-module. Show that the following conditions are equivalent:

\begin{enumerate}
\def\labelenumi{(\roman{enumi})}
\item
  the canonical homomorphism \(\alpha_{\mathrm{M}}: \mathrm{M} \to \mathrm{D}(\mathrm{D}(\mathrm{M}))\) is bijective;
\item
  there exists an exact sequence \(0 \to M' \to M \to M'' \to 0\) , where \(M'\) is of finite type and \(M''\) Artinian.
\end{enumerate}

(If M satisfies (i), show that its socle is of finite dimension over \(\kappa_{A}\) , and that the same is true of every quotient of M; construct a finitely generated submodule N of M such that M/mN is an essential extension of its socle, and use exerc. 1).

\begin{enumerate}
\def\labelenumi{\arabic{enumi})}
\setcounter{enumi}{2}
\tightlist
\item
  Let \(k\) a field, S the ring \(k[\mathrm{T}_{1},\ldots ,\mathrm{T}_{d}]\) , \(\mathfrak{m}\) the ideal \((\mathrm{T}_1,\dots ,\mathrm{T}_d)\) . One denotes by \(\mathrm{S^{*gr}}\) the graded dual of S (no. 6), and \((\mathbf{T}^{-\alpha})_{\alpha \in \mathbf{N}^d}\) the dual basis of \((\mathbf{T}^{\alpha})_{\alpha \in \mathbf{N}^d}\) (denoted \((u_{\alpha})_{\alpha \in \mathbf{N}^d}\) in loc. cit.).
\end{enumerate}

\begin{enumerate}
\def\labelenumi{\alph{enumi})}
\item
  To every finitely generated sub-S-module M of S*gr one associates the ideal Ann(M) of S. Show that this defines a bijective correspondence between the finitely generated submodules of S*gr and the ideals a of S contained in m and such that S/a is of finite length; the inverse bijection associates to the ideal a the submodule M of S*gr consisting of the elements annihilated by a (observe that M is a Matlis S/a-module).
\item
  For S/a to be a Gorenstein ring, it is necessary and sufficient that the S-module M be cyclic.
\item
  Let \(f\) the element \(\sum \mathrm{T}_i^{-2}\) of \(\mathrm{S}^{*\mathrm{gr}}\) ; prove that the ideal \(\mathfrak{a}\) corresponding to the submodule \(A f\) of \(\mathrm{S}^{*\mathrm{gr}}\) is generated by the quadratic forms \(\mathrm{T}_i\mathrm{T}_j\) for \(i < j\) and \(\mathrm{T}_i^2 - \mathrm{T}_1^2\) for \(i > 1\) If \(d \geqslant 3\) the ideal \(\mathfrak{a}\) is not completely secant; the Gorenstein ring \(A / \mathfrak{a}\) is not a complete intersection ring (cf. § 5, exerc. 5).
\end{enumerate}

\begin{enumerate}
\def\labelenumi{\arabic{enumi})}
\setcounter{enumi}{3}
\item
  Let A be a Noetherian ring, M a finitely generated A-module. For every \(\mathfrak{p} \in \operatorname{Spec}(A)\) , we set \(\boldsymbol{\mu}^{i}(\mathfrak{p}, M) = \dim_{\kappa(\mathfrak{p})} \operatorname{Ext}_{A_{\mathfrak{p}}}^{i}(\kappa(\mathfrak{p}), M_{\mathfrak{p}})\) ; we denote \((I(\mathfrak{p}), e_{\mathfrak{p}})\) an injective envelope of the A-module \(A/\mathfrak{p}\) . Let I be a minimal injective resolution of M (which means that \(I^{n}\) is an injective envelope of \(Z^{n}(I)\) for every \(n \geqslant 0\) , cf. A, X, p. 182, exerc. 13). Prove that the A-module \(I^{p}\) is isomorphic to \(\bigoplus_{\mathfrak{p}} I(\mathfrak{p})^{\mu^{p}(\mathfrak{p}, M)}\) .
\item
  Let A be a Noetherian local ring, I a Matlis A-module. For every A-module M one denotes by D(M) the Matlis dual \(\mathrm{Hom}_{\mathrm{A}}(\mathrm{M},\mathrm{I})\) . Recall that one denotes by \(\mathrm{dpl}_{\mathrm{A}}(\mathrm{M})\) the infimum of the lengths of flat resolutions of M (A, X, p. 202, exerc. 7). Prove the equalities \(\mathrm{di}_{\mathrm{A}}(\mathrm{D}(\mathrm{M})) = \mathrm{dpl}_{\mathrm{A}}(\mathrm{M})\) , \(\mathrm{di}_{\mathrm{A}}(\mathrm{M}) = \mathrm{dpl}_{\mathrm{A}}(\mathrm{D}(\mathrm{M}))\) .
\item
  Let A be a Noetherian ring, I and J injective A-modules. Prove that the A-module \(\mathrm{Hom}_{\mathrm{A}}(\mathrm{I},\mathrm{J})\) is flat (cf. prop. 6, b)).
\item
  Let A and B be rings not necessarily commutative, J a (A, B)-bimodule, P a flat left B-module. Assume the ring A is left Noetherian. Prove that if J is an injective A-module, then J ⊗B P is likewise.
\end{enumerate}

¶ 8) Let A be a Noetherian local ring.

\begin{enumerate}
\def\labelenumi{\alph{enumi})}
\tightlist
\item
  Let M and N be A-modules such that \(\mathrm{dpl}_{\mathrm{A}}(\mathrm{M}) < +\infty\) and \(\operatorname{Tor}_{i}^{\mathrm{A}}(\mathrm{M}, \mathrm{N}) = 0\) for i \textgreater{} 0. Define for every integer n a canonical isomorphism
\end{enumerate}

\[
\operatorname{Ext} _ {\mathrm{A}} ^ {n} \left(\kappa_ {\mathrm{A}}, \mathrm{M} \otimes_ {\mathrm{A}} \mathrm{N}\right) \longrightarrow \bigoplus_ {p \geqslant n} \left(\operatorname{Tor} _ {p - n} ^ {\mathrm{A}} \left(\kappa_ {\Lambda}, \mathrm{M}\right) \otimes_ {k} \operatorname{Ext} _ {\mathrm{A}} ^ {p} \left(\kappa_ {\mathrm{A}}, \mathrm{N}\right)\right),
\]

and prove that one has \(\mathrm{di}_{\mathrm{A}}(\mathrm{M}\otimes_{\mathrm{A}}\mathrm{N})\leqslant \mathrm{di}_{\mathrm{A}}(\mathrm{N})\) (adapt the proof of exerc. 6 of § 3 with the aid of exerc. 7).

\begin{enumerate}
\def\labelenumi{\alph{enumi})}
\setcounter{enumi}{1}
\tightlist
\item
  Let N, P be A-modules such that \(\mathrm{di}_{\mathrm{A}}(\mathrm{P}) < +\infty\) and \(\mathrm{Ext}_{\mathrm{A}}^{i}(\mathrm{N},\mathrm{P}) = 0\) for i \textgreater{} 0. Define for every integer n a canonical isomorphism
\end{enumerate}

\[
\operatorname{Tor} _ {n} ^ {\mathrm{A}} \left(\kappa_ {\mathrm{A}}, \operatorname{Hom} _ {\mathrm{A}} (\mathrm{N}, \mathrm{P})\right) \longrightarrow \bigoplus_ {p \geqslant n} \operatorname{Hom} _ {\kappa_ {\mathrm{A}}} \left(\operatorname{Ext} _ {\mathrm{A}} ^ {p - n} \left(\kappa_ {\mathrm{A}}, \mathrm{P}\right), \operatorname{Ext} _ {\mathrm{A}} ^ {p} \left(\kappa_ {\mathrm{A}}, \mathrm{N}\right)\right)
\]

(apply Matlis duality).

\begin{enumerate}
\def\labelenumi{\alph{enumi})}
\setcounter{enumi}{2}
\item
  Prove that if there exists an A-module M whose injective dimension and flat dimension are finite, then A is a Gorenstein ring (take N = A in a)).
\item
  Conversely, if A is a Gorenstein ring, prove that there is an identity between modules of finite injective dimension and modules of finite flat dimension (prove with the aid of a) that a module of finite flat dimension is of finite injective dimension, and obtain the converse by Matlis duality).
\item
  Deduce from d) that the following conditions are equivalent:
\end{enumerate}

\begin{enumerate}
\def\labelenumi{(\roman{enumi})}
\item
  A is a Gorenstein ring;
\item
  the finitely generated A-modules of finite injective dimension coincide with the finitely generated A-modules of finite projective dimension;
\item
  there exists a finitely generated A-module whose injective dimension and projective dimension are finite.
\end{enumerate}

\begin{enumerate}
\def\labelenumi{\arabic{enumi})}
\setcounter{enumi}{8}
\tightlist
\item
  Let A be a Noetherian local ring satisfying condition (GM) (§ 3, exerc. 13).
\end{enumerate}

\begin{enumerate}
\def\labelenumi{\alph{enumi})}
\item
  Let I be a complex of injective A-modules, of length \(\ell\) , whose homology is nonzero and of finite type. Prove the inequality \(\ell \geqslant \dim(A)\) (reduce to loc. cit. a) by Matlis duality).
\item
  If there exists a finitely generated A-module M of finite injective dimension, then A is a Macaulay ring (apply a) to an injective resolution of M).
\item
  For A to be a Gorenstein ring, it is necessary and sufficient that it possess a cyclic module of finite injective dimension (use Exercise 8, b)).
\end{enumerate}

\begin{enumerate}
\def\labelenumi{\arabic{enumi})}
\setcounter{enumi}{9}
\tightlist
\item
  Let A be a Noetherian ring, and let M, P, T be A-modules; suppose that M is finitely generated, T has finite injective dimension, and that the modules \(\operatorname{Ext}_{\mathrm{A}}^{i}(\mathrm{M},\mathrm{P})\) and \(\operatorname{Ext}_{\mathrm{A}}^{i}(\mathrm{P},\mathrm{T})\) are zero for i \textgreater{} 0. Prove that the canonical homomorphism \(\mathrm{M}\otimes_{\mathrm{A}}\mathrm{Hom}_{\mathrm{A}}(\mathrm{P},\mathrm{T})\to\mathrm{Hom}_{\mathrm{A}}(\mathrm{Hom}_{\mathrm{A}}(\mathrm{M},\mathrm{P}),\mathrm{T})\) which maps \(m\otimes h\) onto the homomorphism \(u\mapsto h(u(m))\) is an isomorphism (argue as in the proof of Prop. 6, replacing J by an injective resolution of T).
\end{enumerate}

